\chapter{Comparing Groups and Categorical Data}
\label{ch:comparing-groups}

Chapter 9 built the machinery of a test: a null hypothesis, a test statistic,
its null distribution, and a p-value. It applied that machinery to means with
the $z$- and $t$-tests. Real questions come in more shapes. Did the new
checkout page convert better? Is this die fair? Does conversion depend on
device? Do three server configurations have the same mean latency? What if
the data are skewed and $n=6$? This chapter gives one tool for each shape and
ends with a flowchart for choosing among them.

Every test here follows the same recipe:
\begin{enumerate}
  \item Compute what you would \emph{expect} to see if nothing were going on.
  \item Measure how far the data sit from that expectation, in units of noise.
  \item Compare the distance to its null distribution to get a p-value.
\end{enumerate}

\section{Two proportions: the A/B conversion test}
\label{sec:cg-two-prop}

\subsection{A worked example}

An online shop shows its old checkout page (control, A) to $10{,}000$ randomly
chosen users and a new page (treatment, B) to another $10{,}000$.

\begin{center}
\begin{tabular}{lrrr}
\toprule
 & Users & Conversions & Rate \\
\midrule
A (control) & $10{,}000$ & $1{,}000$ & $\hat p_A = 0.100$ \\
B (treatment) & $10{,}000$ & $1{,}100$ & $\hat p_B = 0.110$ \\
\bottomrule
\end{tabular}
\end{center}

The observed lift is $\hat p_B - \hat p_A = 0.010$, one percentage point. Is
that a real effect, or could two identical pages produce a gap this big by
chance?

Each user converts or not, so each arm is a sample of Bernoulli trials. A
sample proportion is a sample mean of 0s and 1s, so Chapter 6 gives its
standard error: $\SE(\hat p) = \sqrt{p(1-p)/n}$. The two arms are independent,
so variances add:
\[
  \Var(\hat p_B - \hat p_A) = \frac{p_A(1-p_A)}{n_A} + \frac{p_B(1-p_B)}{n_B}.
\]

Under the null hypothesis $H_0\colon p_A = p_B = p$, both arms share one rate.
The best estimate of that common rate pools all users:
\[
  \hat p = \frac{1000 + 1100}{10000 + 10000} = \frac{2100}{20000} = 0.105.
\]
Plug it into the variance formula. The subscript $0$ marks a standard error
computed as if $H_0$ were true:
\[
  \SE_0 = \sqrt{\hat p(1-\hat p)\left(\frac{1}{n_A}+\frac{1}{n_B}\right)}
        = \sqrt{0.105 \times 0.895 \times \frac{2}{10000}}
        = \sqrt{1.8795\times 10^{-5}} \approx 0.004335.
\]
The test statistic counts how many null standard errors the observed gap is:
\[
  z = \frac{\hat p_B - \hat p_A}{\SE_0} = \frac{0.010}{0.004335} \approx 2.31.
\]
By the CLT the statistic is close to $\Normal(0,1)$ under $H_0$. The two-sided
p-value is $2\,\P(Z > 2.31) \approx 0.021$. At $\alpha = 0.05$ we reject $H_0$.

\begin{recoverybox}
Keep the same rates ($10\%$ and $11\%$) but use only $1000$
users per arm. Does $z$ grow or shrink, and by how much? (Answer: the SE grows
by $\sqrt{10}$, so $z$ falls to about $0.73$ and $p \approx 0.47$. Same lift, no
evidence.)
\end{recoverybox}

\begin{figure}[H]
\centering
\begin{tikzpicture}
\begin{axis}[width=11cm,height=5cm,domain=-4:4,samples=120,
  axis lines=left,xlabel={$z$},ylabel={density},ymin=0,ymax=0.45,
  xtick={-2.31,-1,0,1,2.31},xticklabels={$-2.31$,$-1$,$0$,$1$,$2.31$},
  ytick=\empty,clip=false]
  \addplot[draw=none,fill=trapred!45,domain=2.31:4] {exp(-x^2/2)/sqrt(2*pi)} \closedcycle;
  \addplot[draw=none,fill=trapred!45,domain=-4:-2.31] {exp(-x^2/2)/sqrt(2*pi)} \closedcycle;
  \addplot[thick,distblue] {exp(-x^2/2)/sqrt(2*pi)};
  \draw[thick,densityred] (axis cs:2.31,0) -- (axis cs:2.31,0.2);
  \node[anchor=south,font=\small,densityred] at (axis cs:2.31,0.2) {observed $z$};
  \node[anchor=west,font=\small] at (axis cs:2.6,0.08) {$0.0105$};
  \node[anchor=east,font=\small] at (axis cs:-2.6,0.08) {$0.0105$};
  \node[font=\small,distblue] at (axis cs:0,0.43) {null: $\Normal(0,1)$};
\end{axis}
\end{tikzpicture}
\caption{Null distribution of the two-proportion $z$ statistic. The observed
$z=2.31$ leaves $0.0105$ in each tail, so the two-sided p-value is $0.021$: a
gap this large would be rare if the pages were identical.}
\label{fig:cg-abz}
\end{figure}

\subsection{Pooled for the test, unpooled for the interval}

The test assumes $H_0$ is true, so it uses the pooled $\hat p$. A confidence
interval for the lift makes no such assumption, so it uses each arm's own rate:
\[
  \SE = \sqrt{\frac{0.10\times 0.90}{10000} + \frac{0.11\times 0.89}{10000}}
      = \sqrt{9.00\times10^{-6} + 9.79\times10^{-6}} \approx 0.004335.
\]
The 95\% interval is $0.010 \pm 1.96 \times 0.004335 = [0.0015,\ 0.0185]$.
The lift is somewhere between $0.15$ and $1.85$ percentage points. The interval
excludes 0, which agrees with the test. Here the two standard errors match to
four digits because the rates are close; they differ when the rates are far
apart.

\begin{figure}[H]
\centering
\begin{tikzpicture}
\begin{axis}[width=9cm,height=5cm,ymin=0.09,ymax=0.12,xmin=0.4,xmax=2.6,
  xtick={1,2},xticklabels={A (control),B (treatment)},
  ylabel={conversion rate},axis lines=left,
  yticklabel style={/pgf/number format/fixed,/pgf/number format/precision=3}]
  \addplot[only marks,mark=*,mark size=2.5pt,sourceblue,
    error bars/.cd,y dir=both,y explicit]
    coordinates {(1,0.100) +- (0,0.00588) (2,0.110) +- (0,0.00613)};
  \node[anchor=west,font=\small] at (axis cs:1.08,0.100) {$0.100$};
  \node[anchor=west,font=\small] at (axis cs:2.08,0.110) {$0.110$};
\end{axis}
\end{tikzpicture}
\caption{Each arm's conversion rate with its own 95\% interval
($\pm 1.96\sqrt{\hat p(1-\hat p)/n}$). The two bars overlap slightly, yet the
test on the \emph{difference} rejects. Overlapping individual intervals do not
imply a non-significant difference.}
\label{fig:cg-ab-bars}
\end{figure}

\begin{trapbox}
Do not compare two groups by checking whether their separate confidence
intervals overlap. The SE of a difference is
$\sqrt{\SE_A^2 + \SE_B^2}$, which is smaller than $\SE_A + \SE_B$. Two
intervals can overlap while the interval for the difference excludes zero, as in
Figure~\ref{fig:cg-ab-bars}. Always build the interval for the difference itself.
\end{trapbox}

\begin{whybox}
Why is the $z$ approximation good here? Each arm's count of conversions is a
sum of $10{,}000$ independent Bernoulli variables, so the CLT applies to each
proportion. A common rule of thumb: you want at least about 10 successes and 10
failures in each arm. With $1000$ and $9000$ we are far past that. With 4
conversions in 40 users you are not, and Fisher's exact test
(Section~\ref{sec:cg-fisher}) is the safer choice.
\end{whybox}

\begin{lstlisting}[language=Python]
import numpy as np
from scipy import stats
xA, nA, xB, nB = 1000, 10000, 1100, 10000
p = (xA + xB) / (nA + nB)                      # pooled rate 0.105
se0 = np.sqrt(p*(1-p)*(1/nA + 1/nB))
z = (xB/nB - xA/nA) / se0                      # 2.31
print(z, 2*stats.norm.sf(abs(z)))              # p = 0.021
\end{lstlisting}

\section{Chi-square goodness of fit}
\label{sec:cg-gof}

\subsection{Is this die fair?}

You roll a die 120 times. A fair die should give each face $120/6 = 20$
times on average. You observe:

\begin{center}
\begin{tabular}{lrrrrrr|r}
\toprule
Face & 1 & 2 & 3 & 4 & 5 & 6 & Total \\
\midrule
Observed $O$ & 15 & 22 & 18 & 25 & 16 & 24 & 120 \\
Expected $E$ & 20 & 20 & 20 & 20 & 20 & 20 & 120 \\
$O - E$ & $-5$ & $2$ & $-2$ & $5$ & $-4$ & $4$ & 0 \\
$(O-E)^2/E$ & 1.25 & 0.20 & 0.20 & 1.25 & 0.80 & 0.80 & 4.50 \\
\bottomrule
\end{tabular}
\end{center}

\begin{figure}[H]
\centering
\begin{tikzpicture}
\begin{axis}[width=10cm,height=5cm,ybar,bar width=14pt,ymin=0,ymax=30,
  xtick={1,2,3,4,5,6},xlabel={face},ylabel={count},axis lines=left,
  enlarge x limits=0.12,legend columns=2,
  legend style={at={(0.5,1.03)},anchor=south,font=\small,/tikz/every even column/.append style={column sep=8pt}}]
  \addplot[fill=histfill,draw=distblue] coordinates {(1,15) (2,22) (3,18) (4,25) (5,16) (6,24)};
  \addlegendentry{observed}
  \addplot[sharp plot,ultra thick,densityred,dashed,mark=none,forget plot] coordinates {(0.5,20) (6.5,20)};
  \addlegendimage{ultra thick,densityred,dashed}
  \addlegendentry{expected $=20$}
\end{axis}
\end{tikzpicture}
\caption{Observed die counts against the fair expectation. The bars wobble
around 20 by up to 5. The chi-square statistic decides whether that wobble is
more than chance.}
\label{fig:cg-die}
\end{figure}

\begin{definition}[Pearson chi-square statistic]
For $k$ categories with observed counts $O_i$ and expected counts $E_i$ under
$H_0$,
\[
  X^2 = \sum_{i=1}^{k} \frac{(O_i - E_i)^2}{E_i}.
\]
For a goodness-of-fit test with fully specified probabilities, $X^2$ is
approximately $\chi^2_{k-1}$ under $H_0$.
\end{definition}

The subscript $k-1$ is the \emph{degrees of freedom}: the number of counts
that are free to vary once the total is fixed. The chi-square distribution
$\chi^2_m$ is the distribution of a sum of $m$ squared independent standard
normal variables, so it describes ``total squared surprise'' from $m$ free
pieces.

For the die, $X^2 = 4.50$ with $k - 1 = 5$ degrees of freedom. The 95th
percentile of $\chi^2_5$ is $11.07$, and $\P(\chi^2_5 > 4.5) \approx 0.48$.
We have no evidence against fairness. Wobbles of this size happen about half
the time with a fair die.

\begin{whybox}
Why divide by $E_i$? A count's natural wobble grows with its size: a cell
expecting 1000 can easily miss by 30, a cell expecting 10 cannot. Each count
$O_i$ has variance $E_i(1-p_i)$, close to $E_i$, so $(O_i - E_i)/\sqrt{E_i}$ is
roughly a standardised deviation and $X^2$ adds up squared standardised
deviations. A sum of $m$ independent squared standard normals is the
definition of $\chi^2_m$. Why $k-1$ and not $k$? The counts must add to $n$,
so once five deviations are known the sixth is fixed: $\sum (O_i - E_i) = 0$.
The deviations are not independent; that one constraint (together with the
$1-p_i$ factors left out of the denominators) removes exactly one degree of
freedom. (See Probability From the Ground Up
for the chi-square distribution itself.)
\end{whybox}

\subsection{The chi-square family}

The chi-square distribution with $m$ degrees of freedom has mean $m$ and
variance $2m$. Its shape depends only on $m$. Small $m$ gives a right-skewed
curve piled near zero; large $m$ looks more and more normal.

\begin{figure}[H]
\centering
\begin{tikzpicture}
\begin{axis}[width=11.5cm,height=5.5cm,xmin=0,xmax=20,ymin=0,ymax=0.5,
  axis lines=left,xlabel={$x$},ylabel={density},samples=150,
  legend style={at={(0.98,0.98)},anchor=north east,font=\small},
  restrict y to domain=0:0.6]
  \addplot[thick,distblue,domain=0.03:20] {exp(-x/2)/sqrt(2*pi*x)};
  \addlegendentry{$m=1$}
  \addplot[thick,targetgreen,domain=0:20] {0.5*exp(-x/2)};
  \addlegendentry{$m=2$}
  \addplot[thick,quotientorange,domain=0:20] {x*exp(-x/2)/4};
  \addlegendentry{$m=4$}
  \addplot[thick,structurepurple,domain=0:20] {x^2*exp(-x/2)/16};
  \addlegendentry{$m=6$}
  \addplot[thick,densityred,domain=0:20] {x^4*exp(-x/2)/768};
  \addlegendentry{$m=10$}
\end{axis}
\end{tikzpicture}
\caption{Chi-square densities for several degrees of freedom $m$. Each has
mean $m$ (and, for $m \ge 2$, its peak at $m-2$), so the same $X^2$ value means very different
things at different $m$: $X^2 = 8$ is extreme for $m=1$ and ordinary for
$m=10$.}
\label{fig:cg-chisq-family}
\end{figure}

\begin{figure}[H]
\centering
\begin{tikzpicture}
\begin{axis}[width=11cm,height=5cm,xmin=0,xmax=20,ymin=0,ymax=0.18,
  axis lines=left,xlabel={$X^2$},ylabel={density},samples=150,ytick=\empty,
  xtick={0,4.5,11.07,15,20},xticklabels={$0$,$4.5$,$11.07$,$15$,$20$},clip=false]
  \addplot[draw=none,fill=histfill!60,domain=4.5:20] {x^1.5*exp(-x/2)/7.5199} \closedcycle;
  \addplot[draw=none,fill=trapred!50,domain=11.07:20] {x^1.5*exp(-x/2)/7.5199} \closedcycle;
  \addplot[thick,distblue,domain=0.001:20] {x^1.5*exp(-x/2)/7.5199};
  \draw[thick,densityred] (axis cs:4.5,0) -- (axis cs:4.5,0.16);
  \node[anchor=south west,font=\small] at (axis cs:4.6,0.155) {observed $4.5$: $p=0.48$};
  \draw[->] (axis cs:13.5,0.06) -- (axis cs:12.2,0.012);
  \node[anchor=south,font=\small,trapred] at (axis cs:13.5,0.06) {rejection region, 5\%};
\end{axis}
\end{tikzpicture}
\caption{The $\chi^2_5$ null distribution for the die. The blue area to the
right of the observed $4.5$ is the p-value, $0.48$. To reject at 5\% we would
need $X^2 > 11.07$.}
\label{fig:cg-die-tail}
\end{figure}

\begin{trapbox}
Chi-square tests work on \emph{counts}, never on percentages or rates. If you
feed in $12.5\%, 18.3\%, \dots$ instead of the counts, $X^2$ changes with the
units and the p-value is meaningless. Also check that every expected count is
at least about 5. If not, merge sparse categories or use an exact test.
\end{trapbox}

\begin{historybox}
Karl Pearson, working in London in 1900, wanted to test whether data fitted a
proposed curve: did a set of dice throws, or measured traits, match the
frequencies a theory predicted? Earlier methods judged fit by eye. Pearson
introduced $X^2 = \sum (O-E)^2/E$ and worked out its large-sample distribution,
the chi-square distribution. He got the degrees of freedom wrong for tables
whose expected counts are estimated from the data; Ronald Fisher corrected this
in 1922, giving the $(r-1)(c-1)$ rule. Fisher went on, at the Rothamsted
agricultural station in the 1920s, to compare crop yields across many
fertiliser treatments at once. He split total variation into between-plot and
within-plot pieces, which became the analysis of variance; the $F$ distribution
is named in his honour. Together these gave the first general tools for
comparing more than two groups.
\end{historybox}

\section{Chi-square test of independence}
\label{sec:cg-indep}

\subsection{Does conversion depend on device?}

A site records 1000 visits by device and whether the visit converted.

\begin{center}
\begin{tabular}{lrr|r}
\toprule
 & Converted & Not converted & Row total \\
\midrule
Desktop & 60 & 340 & 400 \\
Mobile & 50 & 450 & 500 \\
Tablet & 10 & 90 & 100 \\
\midrule
Column total & 120 & 880 & 1000 \\
\bottomrule
\end{tabular}
\end{center}

Desktop converts at $60/400 = 15\%$; mobile and tablet at $10\%$. The null
hypothesis is that device and conversion are \emph{independent}: every device
has the same conversion rate, which must then be the overall
$120/1000 = 12\%$.

\textbf{Expected counts.} If independent, desktop should have $12\%$ of its 400
visits converting: $0.12 \times 400 = 48$. In general,
\[
  E_{ij} = \frac{(\text{row } i \text{ total}) \times (\text{column } j \text{ total})}{\text{grand total}}
         = \frac{R_i\, C_j}{N},
\]
where $R_i$ is the row total, $C_j$ the column total and $N$ the grand total.
For desktop and converted: $400 \times 120 / 1000 = 48$.

\begin{recoverybox}
Why must the expected table have the same row and column
totals as the observed table? (Because $\sum_j E_{ij} = R_i \sum_j C_j / N =
R_i$. Independence only reshuffles counts inside the margins.)
\end{recoverybox}

The full expected
table and each cell's contribution $(O-E)^2/E$:

\begin{center}
\begin{tabular}{lrr|rr}
\toprule
 & $E$: Conv. & $E$: Not & $\frac{(O-E)^2}{E}$: Conv. & $\frac{(O-E)^2}{E}$: Not \\
\midrule
Desktop & 48 & 352 & $144/48 = 3.000$ & $144/352 = 0.409$ \\
Mobile & 60 & 440 & $100/60 = 1.667$ & $100/440 = 0.227$ \\
Tablet & 12 & 88 & $4/12 = 0.333$ & $4/88 = 0.045$ \\
\bottomrule
\end{tabular}
\end{center}

Adding the six contributions gives $X^2 \approx 5.68$.

\textbf{Degrees of freedom.} For an $r\times c$ table, $\text{df} = (r-1)(c-1)$.
Here $(3-1)(2-1) = 2$. The reason: once the row and column totals are fixed,
you can fill only $(r-1)(c-1)$ cells freely; the rest are forced by the totals.

\textbf{Conclusion.} The 95th percentile of $\chi^2_2$ is $5.99$, and
$\P(\chi^2_2 > 5.68) \approx 0.058$. The data hint that desktop converts
better, but at $\alpha = 0.05$ the evidence falls just short. The cell
contributions show where the signal is: desktop-converted contributes $3.0$ of
the $5.68$.

\begin{figure}[H]
\centering
\begin{tikzpicture}
\begin{axis}[width=10cm,height=5.2cm,ybar,bar width=16pt,ymin=0,ymax=110,
  symbolic x coords={Desktop,Mobile,Tablet},xtick=data,
  ylabel={conversions},axis lines=left,enlarge x limits=0.25,
  legend style={at={(0.98,0.98)},anchor=north east,font=\small},
  nodes near coords,nodes near coords style={font=\scriptsize}]
  \addplot[fill=histfill,draw=distblue] coordinates {(Desktop,60) (Mobile,50) (Tablet,10)};
  \addlegendentry{observed $O$}
  \addplot[fill=quotientorange!50,draw=quotientorange] coordinates {(Desktop,48) (Mobile,60) (Tablet,12)};
  \addlegendentry{expected $E$ if independent}
\end{axis}
\end{tikzpicture}
\caption{Observed versus expected conversions per device. Desktop has 12 more
conversions than independence predicts and mobile has 10 fewer. Desktop's cell
contributes most to $X^2$: its gap is the largest and its expected count (48)
is smaller than mobile's (60), so $12^2/48 = 3.0$ beats $10^2/60 = 1.7$.}
\label{fig:cg-device}
\end{figure}

\begin{connectionbox}
A $2\times 2$ table is the A/B test in disguise. Put the checkout data in a
table (rows A and B; columns converted, not converted) and run the chi-square
test without continuity correction: $X^2 = 5.32$. That equals
$z^2 = 2.3066^2 = 5.32$, and the p-values are identical ($0.021$). The
two-proportion $z$-test and the $2\times 2$ chi-square test are the same test.
The $z$ form has one advantage: it tells you the \emph{direction} and gives a
confidence interval for the lift.
\end{connectionbox}

\begin{trapbox}
Independence tests detect association, not causation. If desktop users are
mostly office workers buying on lunch break, device may be a proxy for intent.
Only randomisation (Chapter 5) licenses a causal reading. The checkout A/B test
is causal; the device table is not.
\end{trapbox}

\section{Fisher's exact test}
\label{sec:cg-fisher}

When counts are small, the chi-square approximation breaks down. Fisher's exact
test computes the p-value directly from the hypergeometric distribution, with no
large-sample approximation.

A pilot gives a new onboarding flow to 10 users and the old one to 10. Results:

\begin{center}
\begin{tabular}{lrr|r}
\toprule
 & Retained & Churned & Total \\
\midrule
New flow & 7 & 3 & 10 \\
Old flow & 2 & 8 & 10 \\
\midrule
Total & 9 & 11 & 20 \\
\bottomrule
\end{tabular}
\end{center}

\textbf{The idea.} Fix the margins: 9 retained users among 20, 10 of whom got the
new flow. If the flow has no effect, which 10 users received it is a random
draw, and the number of retained users in the new-flow group $X$ follows a
hypergeometric distribution: draw 10 from an urn of 20 containing 9 ``retained''
balls. Then
\[
  \P(X = k) = \frac{\binom{9}{k}\binom{11}{10-k}}{\binom{20}{10}}.
\]
Observing $X = 7$ or more has probability
$0.0322 + 0.0027 + 0.0001 = 0.035$. The distribution is symmetric here, so the
two-sided p-value is about $0.070$.

\begin{figure}[H]
\centering
\begin{tikzpicture}
\begin{axis}[width=10.5cm,height=5cm,ybar,bar width=12pt,ymin=0,ymax=0.36,
  xtick={0,...,9},xlabel={retained users in new-flow group, $X$},
  ylabel={$\P(X=k)$},axis lines=left,enlarge x limits=0.06,clip=false]
  \addplot[fill=histfill,draw=distblue,bar shift=0pt] coordinates
    {(3,0.1500) (4,0.3151) (5,0.3151) (6,0.1500)};
  \addplot[fill=trapred!60,draw=trapred,bar shift=0pt] coordinates
    {(0,0.0001) (1,0.0027) (2,0.0322) (7,0.0322) (8,0.0027) (9,0.0001)};
  \draw[->] (axis cs:8.2,0.14) -- (axis cs:7.1,0.045);
  \node[anchor=south,font=\small] at (axis cs:8.2,0.14) {observed $X=7$};
\end{axis}
\end{tikzpicture}
\caption{Null distribution of $X$ for Fisher's exact test with margins 9/11
and 10/10. Values $X \ge 7$ have total probability $0.035$. Adding the equally
extreme lower tail $X \le 2$ (also red) gives a two-sided p-value of $0.070$.}
\label{fig:cg-fisher}
\end{figure}

The chi-square test on the same table gives $X^2 = 5.05$ and $p = 0.025$. That
is too optimistic: two of the expected counts are $4.5$, below the rule of
thumb. With small counts, trust Fisher. With large counts the two agree and
the chi-square test is faster.

\section{One-way ANOVA: comparing several means}
\label{sec:cg-anova}

\subsection{Three server configurations}

You measure p50 latency (ms) on 5 runs each for three server configurations:

\begin{center}
\begin{tabular}{lccccc|c}
\toprule
Config & \multicolumn{5}{c|}{Runs} & Mean \\
\midrule
A & 12 & 14 & 11 & 13 & 15 & 13 \\
B & 16 & 18 & 15 & 17 & 19 & 17 \\
C & 13 & 15 & 14 & 12 & 16 & 14 \\
\bottomrule
\end{tabular}
\end{center}

The grand mean of all 15 runs is $220/15 \approx 14.667$. Are the three
population means equal? You could run three pairwise $t$-tests, but each has a
5\% false alarm rate, so the chance of at least one false alarm is about
$1 - 0.95^3 \approx 14\%$ (Chapter 11 treats this properly). ANOVA asks
one question: $H_0\colon \mu_A = \mu_B = \mu_C$. ANOVA stands for
\emph{analysis of variance}: it decides whether means differ by comparing two
kinds of spread, the spread of the group means and the spread inside groups.

\begin{figure}[H]
\centering
\begin{tikzpicture}
\begin{axis}[name=left,width=6.3cm,height=6cm,xmin=0.4,xmax=3.6,ymin=4,ymax=25,
  xtick={1,2,3},xticklabels={A,B,C},ylabel={latency (ms)},axis lines=left,
  title={\small small within-group spread},clip=false]
  \addplot[only marks,mark=*,mark size=1.8pt,sourceblue] coordinates
   {(1,12) (1,14) (1,11) (1,13) (1,15) (2,16) (2,18) (2,15) (2,17) (2,19) (3,13) (3,15) (3,14) (3,12) (3,16)};
  \draw[very thick,densityred] (axis cs:0.75,13) -- (axis cs:1.25,13);
  \draw[very thick,densityred] (axis cs:1.75,17) -- (axis cs:2.25,17);
  \draw[very thick,densityred] (axis cs:2.75,14) -- (axis cs:3.25,14);
  \draw[dashed,residgray] (axis cs:0.4,14.667) -- (axis cs:3.6,14.667);
  \node[font=\scriptsize,anchor=north] at (axis cs:2,6) {$F = 8.67$, $p = 0.005$};
\end{axis}
\begin{axis}[at={(left.east)},anchor=west,xshift=1.2cm,width=6.3cm,height=6cm,
  xmin=0.4,xmax=3.6,ymin=4,ymax=25,
  xtick={1,2,3},xticklabels={A$'$,B$'$,C$'$},axis lines=left,
  title={\small large within-group spread},clip=false]
  \addplot[only marks,mark=*,mark size=1.8pt,sourceblue] coordinates
   {(1,10) (1,16) (1,7) (1,13) (1,19) (2,14) (2,20) (2,11) (2,17) (2,23) (3,11) (3,17) (3,8) (3,14) (3,20)};
  \draw[very thick,densityred] (axis cs:0.75,13) -- (axis cs:1.25,13);
  \draw[very thick,densityred] (axis cs:1.75,17) -- (axis cs:2.25,17);
  \draw[very thick,densityred] (axis cs:2.75,14) -- (axis cs:3.25,14);
  \draw[dashed,residgray] (axis cs:0.4,14.667) -- (axis cs:3.6,14.667);
  \node[font=\scriptsize,anchor=north] at (axis cs:2,6) {$F = 0.96$, $p = 0.41$};
\end{axis}
\end{tikzpicture}
\caption{Both panels have the same group means (red bars: 13, 17, 14) and the
same grand mean (dashed). On the left the groups barely overlap, so the gaps
between means stand out. On the right each point is three times farther from
its group mean, and the same gaps are lost in the noise. ANOVA measures exactly
this ratio.}
\label{fig:cg-anova-dots}
\end{figure}

\subsection{Splitting the variation}

Take any single observation $x_{ij}$ (run $j$ in group $i$). Its distance from
the grand mean $\bar x$ splits into two pieces:
\[
  \underbrace{x_{ij} - \bar x}_{\text{total}}
  = \underbrace{(\bar x_i - \bar x)}_{\text{between: group mean vs grand mean}}
  + \underbrace{(x_{ij} - \bar x_i)}_{\text{within: point vs its group mean}}.
\]

\begin{figure}[H]
\centering
\begin{tikzpicture}[x=1cm,y=0.55cm]
  \draw[->] (0,9.5) -- (0,20.5) node[above,font=\small] {ms};
  \foreach \v in {10,12,...,20} \draw (-0.1,\v) -- (0.1,\v) node[left=3pt,font=\scriptsize] {\v};
  \draw[dashed,residgray] (0,14.667) -- (10,14.667);
  \node[anchor=north east,font=\small,residgray] at (10,14.5) {grand mean $14.67$};
  \draw[very thick,densityred] (2.5,17) -- (4.5,17) node[right,font=\small] {group B mean $17$};
  \fill[sourceblue] (3.5,19) circle (3pt) node[right=4pt,font=\small] {run $x_{B5} = 19$};
  \draw[very thick,quotientorange,-{Stealth}] (2.2,14.667) -- (2.2,17);
  \node[anchor=east,font=\small,quotientorange,align=right] at (2.1,15.8) {between\\$+2.33$};
  \draw[very thick,sourceblue,-{Stealth}] (2.2,17) -- (2.2,19);
  \node[anchor=east,font=\small,sourceblue,align=right] at (2.1,18.1) {within\\$+2$};
  \draw[thick,structurepurple,-{Stealth}] (9.2,14.667) -- (9.2,19);
  \node[anchor=west,font=\small,structurepurple] at (9.3,17.3) {total $+4.33$};
\end{tikzpicture}
\caption{The deviation of one observation from the grand mean is the sum of a
between-group piece (orange, how far its group mean sits from the grand mean)
and a within-group piece (blue, how far the point sits from its own group
mean).}
\label{fig:cg-decomp}
\end{figure}

Square both sides and sum over all observations. The cross term vanishes,
because within each group the deviations $x_{ij} - \bar x_i$ sum to zero. What
remains is the ANOVA identity:
\[
  \underbrace{\sum_{i,j}(x_{ij}-\bar x)^2}_{\text{SST}}
  = \underbrace{\sum_{i} n_i(\bar x_i - \bar x)^2}_{\text{SSB}}
  + \underbrace{\sum_{i,j}(x_{ij}-\bar x_i)^2}_{\text{SSW}}.
\]

For the latency data:
\begin{align*}
  \text{SSB} &= 5\left[(13-14.667)^2 + (17-14.667)^2 + (14-14.667)^2\right]
             = 5(2.778 + 5.444 + 0.444) = 43.33,\\
  \text{SSW} &= \underbrace{(1+1+4+0+4)}_{A} + \underbrace{(1+1+4+0+4)}_{B}
             + \underbrace{(1+1+0+4+4)}_{C} = 30,\\
  \text{SST} &= 43.33 + 30 = 73.33.
\end{align*}

\subsection{The F statistic}

Sums of squares grow with the number of terms, so divide each by its degrees of
freedom to get a \emph{mean square}. With $k$ groups and $N$ observations in
total:
\[
  \text{MSB} = \frac{\text{SSB}}{k-1}, \qquad
  \text{MSW} = \frac{\text{SSW}}{N-k}, \qquad
  F = \frac{\text{MSB}}{\text{MSW}}.
\]
Here $\text{MSB} = 43.33/2 = 21.67$ and $\text{MSW} = 30/12 = 2.5$, so
$F = 21.67 / 2.5 = 8.67$.

\begin{center}
\begin{tabular}{lrrrr}
\toprule
Source & SS & df & MS & $F$ \\
\midrule
Between groups & 43.33 & 2 & 21.67 & 8.67 \\
Within groups & 30.00 & 12 & 2.50 & \\
\midrule
Total & 73.33 & 14 & & \\
\bottomrule
\end{tabular}
\end{center}

\begin{whybox}
Why is this ratio the right statistic? MSW is a pooled sample variance: it
estimates the noise variance $\sigma^2$ whether or not $H_0$ holds, since it
only looks inside groups. MSB estimates $\sigma^2$ too \emph{if} $H_0$ is true:
the group means are then sample means of size $n$ from one population, so they
vary with variance $\sigma^2/n$, and multiplying their spread by $n$ recovers
$\sigma^2$. If the true means differ, MSB picks up that extra spread and grows.
So $F \approx 1$ under $H_0$ and $F \gg 1$ when means differ. The degrees of
freedom follow the usual rule: $k$ group means minus 1 for the grand mean, and
$N$ observations minus $k$ estimated group means.
\end{whybox}

\subsection{The F distribution}

If the groups are independent, normal, and share one variance, then under
$H_0$ the statistic follows an $F$ distribution with $(k-1, N-k)$ degrees of
freedom. Formally, $F_{d_1,d_2}$ is the distribution of
$(U_1/d_1)/(U_2/d_2)$ for independent $U_1\sim\chi^2_{d_1}$ and
$U_2\sim\chi^2_{d_2}$: a ratio of two variance estimates.

For our data, $F_{2,12}$ has 95th percentile $3.89$, and
$\P(F_{2,12} > 8.67) \approx 0.0047$. We reject equal means.

\begin{recoverybox}
What value of $F$ do you expect if the null hypothesis is
true? Could $F$ ever be below 1? (About 1, since both mean squares estimate
$\sigma^2$. Yes: if group means happen to be closer together than noise alone
predicts, MSB $<$ MSW.)
\end{recoverybox}

\begin{figure}[H]
\centering
\begin{tikzpicture}
\begin{axis}[width=11.5cm,height=5.5cm,xmin=0,xmax=10,ymin=0,ymax=1.05,
  axis lines=left,xlabel={$F$},ylabel={density},samples=150,
  legend style={at={(0.98,0.98)},anchor=north east,font=\small},clip=false]
  \addplot[draw=none,fill=trapred!40,domain=3.89:10,forget plot] {(1+x/6)^(-7)} \closedcycle;
  \addplot[thick,distblue,domain=0:10] {(1+x/6)^(-7)};
  \addlegendentry{$F_{2,12}$}
  \addplot[thick,quotientorange,dashed] coordinates {(0.00,0.0000) (0.10,0.2059) (0.20,0.4309) (0.30,0.5899) (0.40,0.6813) (0.50,0.7190) (0.60,0.7181) (0.70,0.6916) (0.80,0.6494) (0.90,0.5989) (1.00,0.5449) (1.10,0.4908) (1.20,0.4387) (1.30,0.3898) (1.40,0.3449) (1.50,0.3041) (1.60,0.2675) (1.70,0.2348) (1.80,0.2058) (1.90,0.1803) (2.00,0.1578) (2.20,0.1208) (2.40,0.0926) (2.60,0.0711) (2.80,0.0547) (3.00,0.0422) (3.30,0.0288) (3.60,0.0198) (4.00,0.0122) (4.50,0.0068) (5.00,0.0039) (5.50,0.0023) (6.00,0.0014) (7.00,0.0005) (8.00,0.0002) (10.0,0.0000)};
  \addlegendentry{$F_{5,20}$}
  \draw[thick,densityred] (axis cs:8.67,0) -- (axis cs:8.67,0.3);
  \node[anchor=south,font=\small,densityred] at (axis cs:8.67,0.3) {observed $8.67$};
  \node[anchor=south,font=\small,trapred] at (axis cs:5.2,0.12) {5\% beyond $3.89$};
  \draw[->,trapred] (axis cs:5.2,0.12) -- (axis cs:4.3,0.03);
\end{axis}
\end{tikzpicture}
\caption{Two $F$ densities. Both are right-skewed with mean near 1, since
under $H_0$ both mean squares estimate the same $\sigma^2$. For our test,
$F_{2,12}$ puts 5\% of its mass beyond $3.89$; the observed $8.67$ is far in
the tail ($p = 0.0047$).}
\label{fig:cg-fdist}
\end{figure}

\subsection{ANOVA with two groups is the t-test}

Run ANOVA on groups A and C only. The means are 13 and 14; each group has
within sum of squares 10. Then $\text{SSB} = 5(0.5^2) + 5(0.5^2) = 2.5$ with
1 df, and $\text{MSW} = 20/8 = 2.5$, so $F = 1$. The pooled two-sample $t$
statistic (Chapter 9) is
\[
  t = \frac{14 - 13}{\sqrt{2.5\,(1/5 + 1/5)}} = \frac{1}{1} = 1.
\]
And $t^2 = 1 = F$. For A versus B, $t = 4$ and $F = 16$, both with
$p = 0.0039$. This holds in general: with $k=2$, $F = t^2$ and $F_{1,d}$ is the
distribution of the square of a $t_d$ variable. One-way ANOVA is the pooled
$t$-test extended to many groups.

\begin{trapbox}
ANOVA assumes equal variances across groups. If one group is much noisier
(a common situation with latencies), use Welch's ANOVA, the many-group cousin
of Welch's $t$-test, or a rank-based test (Section~\ref{sec:cg-nonpar}).
Moderate non-normality is fine with a few dozen points per group, by the CLT
acting on the group means.
\end{trapbox}

\subsection{Post-hoc comparisons}

A significant $F$ says ``not all means are equal''. It does not say which ones
differ. \emph{Post-hoc} tests (Latin for ``after this'': tests run after the
overall $F$ test) compare pairs while keeping the overall false alarm rate
at $\alpha$.

\begin{itemize}
  \item \textbf{Tukey's HSD} (``honestly significant difference'') compares
  every pair, using a critical value from the distribution of the
  \emph{largest} gap among all the group means, so the number of pairs is
  built in. For the latency data: B vs A,
  difference $4$, $p = 0.005$; B vs C, difference $3$, $p = 0.028$; C vs A,
  difference $1$, $p = 0.59$. Config B is slower than both others; A and C are
  indistinguishable.
  \item \textbf{Bonferroni} runs ordinary pairwise tests at level
  $\alpha/m$ for $m$ comparisons. It is simple and conservative.
\end{itemize}

Chapter 11 explains multiple-comparison corrections in depth.

\begin{lstlisting}[language=Python]
from scipy import stats
A = [12,14,11,13,15]; B = [16,18,15,17,19]; C = [13,15,14,12,16]
print(stats.f_oneway(A, B, C))     # F = 8.67, p = 0.0047
print(stats.tukey_hsd(A, B, C))    # pairwise CIs and p-values
print(stats.kruskal(A, B, C))      # rank-based alternative
\end{lstlisting}

\section{Nonparametric alternatives: tests on ranks}
\label{sec:cg-nonpar}

The $t$-test and ANOVA compare means and lean on normality for small samples.
Latencies, incomes and session lengths are skewed and have outliers. Rank
tests replace each value by its position in the sorted data. A 10-second
outlier becomes ``rank 11'', no more influential than a value just above the
rest.

\subsection{Mann--Whitney / Wilcoxon rank-sum test}

Page load times (ms) for two CDNs:
\[
  \text{A: } 12,\ 15,\ 19,\ 22,\ 30 \qquad
  \text{B: } 18,\ 25,\ 28,\ 35,\ 40,\ 41.
\]
Pool the 11 values, sort, and rank them 1 to 11:

\begin{figure}[H]
\centering
\begin{tikzpicture}[x=0.28cm]
  \draw[->] (9,0) -- (44,0) node[right,font=\small] {ms};
  \foreach \v in {10,20,30,40} \draw (\v,0.1) -- (\v,-0.1) node[below,font=\scriptsize] {\v};
  \foreach \v/\r in {12/1,15/2,19/4,22/5,30/8} {
    \fill[sourceblue] (\v,0.6) circle (4pt);
    \node[font=\scriptsize,sourceblue,above] at (\v,0.75) {\r};
  }
  \foreach \v/\r/\o in {18/3/0,25/6/0,28/7/0,35/9/0,40/10/-0.7,41/11/0.7} {
    \fill[densityred] (\v-0.25,-1.05) rectangle (\v+0.25,-0.65);
    \node[font=\scriptsize,densityred,below] at (\v+\o,-1.1) {\r};
  }
  \node[anchor=east,font=\small,sourceblue] at (8.5,0.6) {A};
  \node[anchor=east,font=\small,densityred] at (8.5,-0.85) {B};
\end{tikzpicture}
\caption{Pooled ranks for the two CDNs (A above the axis, B below; numbers are
ranks). A's ranks $1,2,4,5,8$ cluster at the low end, so A tends to be faster.
The rank-sum test asks whether such clustering is surprising if both CDNs
had the same distribution.}
\label{fig:cg-ranks}
\end{figure}

A's rank sum is $R_A = 1 + 2 + 4 + 5 + 8 = 20$. The Mann--Whitney statistic
subtracts the smallest possible rank sum, $n_A(n_A+1)/2 = 15$:
\[
  U_A = R_A - \frac{n_A(n_A+1)}{2} = 20 - 15 = 5.
\]
$U_A$ has a direct meaning: it counts the pairs (one from A, one from B) in
which the A value is larger. Check: 19 and 22 each beat 18; 30 beats 18, 25,
28. That is $1+1+3 = 5$ of the $5 \times 6 = 30$ pairs. So an A load beats a
B load only $5/30 \approx 17\%$ of the time.

Under $H_0$ (same distribution), every assignment of the 11 ranks to the two
groups is equally likely. That gives an exact null distribution for $U$ with
mean $n_A n_B / 2 = 15$. The exact two-sided p-value is $0.082$. The normal
approximation, with $\SD(U) = \sqrt{n_An_B(n_A+n_B+1)/12} = \sqrt{30} = 5.48$,
gives $z = (5 - 15)/5.48 = -1.83$ and $p = 0.068$. With $n = 5$ and $6$ we
lack evidence at 5\%, though the direction is suggestive.

\begin{intuitionbox}
The Mann--Whitney test estimates $\P(X_A > X_B)$, the chance that a random A
value exceeds a random B value. If both groups come from the same
distribution this probability is $1/2$, and the test asks whether the data are
far from that. It is not a test
of medians unless the two distributions have the same shape. That probability
is often what a product manager actually wants: ``how often is the new page
faster?''
\end{intuitionbox}

\subsection{Wilcoxon signed-rank test}

For \emph{paired} data use differences. Eight servers have latency measured
before and after a kernel patch:

\begin{center}
\begin{tabular}{lrrrrrrrr}
\toprule
Server & 1 & 2 & 3 & 4 & 5 & 6 & 7 & 8 \\
\midrule
Before & 120 & 135 & 128 & 150 & 142 & 138 & 160 & 131 \\
After & 112 & 130 & 129 & 138 & 135 & 131 & 149 & 127 \\
$d$ = Before $-$ After & 8 & 5 & $-1$ & 12 & 7 & 7 & 11 & 4 \\
Rank of $|d|$ & 6 & 3 & 1 & 8 & 4.5 & 4.5 & 7 & 2 \\
\bottomrule
\end{tabular}
\end{center}

Rank the absolute differences (ties share the average rank, so the two 7s get
4.5 each). Then add the ranks of the positive and negative differences
separately:
\[
  W_+ = 6 + 3 + 8 + 4.5 + 4.5 + 7 + 2 = 35, \qquad W_- = 1.
\]
The ranks total $1 + 2 + \dots + 8 = 36$. Under $H_0$ (differences symmetric
around zero), each rank is equally likely to carry either sign, giving
$2^8 = 256$ equally likely sign patterns. Only 2 patterns have $W_- \le 1$
(all positive, or only rank 1 negative). The two-sided p-value is
$2 \times 2/256 = 0.0156$. The patch reduced latency.

\begin{figure}[H]
\centering
\begin{tikzpicture}[x=0.75cm]
  \draw[->] (-2.5,0) -- (13.5,0) node[right,font=\small] {$d$ (ms)};
  \foreach \v in {-2,0,2,4,6,8,10,12} \draw (\v,0.1) -- (\v,-0.1) node[below,font=\scriptsize] {\v};
  \draw[dashed,residgray] (0,-0.2) -- (0,2.2);
  \foreach \v/\r/\h in {8/6/0.5,5/3/0.5,12/8/0.5,7/{4.5}/0.5,7/{4.5}/1.25,11/7/0.5,4/2/0.5} {
    \fill[targetgreen] (\v,\h) circle (3.5pt);
    \node[font=\scriptsize,anchor=west] at (\v+0.15,\h+0.15) {\r};
  }
  \fill[trapred] (-1,0.5) circle (3.5pt);
  \node[font=\scriptsize,anchor=east,trapred] at (-1.15,0.65) {1};
  \node[font=\small,targetgreen,anchor=west] at (3.5,2) {$W_+ = 35$};
  \node[font=\small,trapred,anchor=east] at (-0.3,2) {$W_- = 1$};
\end{tikzpicture}
\caption{Paired differences with the rank of $|d|$ beside each point. Only one
difference is negative, and it has the smallest magnitude (rank 1). If the
patch did nothing, signs would be scattered at random across ranks.}
\label{fig:cg-signedrank}
\end{figure}

\subsection{The sign test}

The crudest paired test throws away magnitudes and keeps only signs. If the
patch did nothing, each difference is positive with probability $1/2$, so the
number of positive differences is $\Bin(8, 1/2)$. We saw 7 positives. The
two-sided p-value is
\[
  2\,\P(\text{Bin}(8,\tfrac12) \ge 7) = 2 \cdot \frac{\binom{8}{7}+\binom{8}{8}}{256}
  = 2\cdot\frac{9}{256} \approx 0.070.
\]

\begin{figure}[H]
\centering
\begin{tikzpicture}
\begin{axis}[width=10cm,height=4.8cm,ybar,bar width=12pt,ymin=0,ymax=0.3,
  xtick={0,...,8},xlabel={number of positive differences},
  ylabel={probability},axis lines=left,enlarge x limits=0.07]
  \addplot[fill=trapred!60,draw=trapred,bar shift=0pt] coordinates
    {(0,0.0039) (1,0.03125) (7,0.03125) (8,0.0039)};
  \addplot[fill=histfill,draw=distblue,bar shift=0pt] coordinates
    {(2,0.1094) (3,0.2188) (4,0.2734) (5,0.2188) (6,0.1094)};
\end{axis}
\end{tikzpicture}
\caption{The $\Bin(8,1/2)$ null distribution of the sign test. Outcomes at
least as extreme as 7 positives (red, both tails) have total probability
$18/256 \approx 0.070$.}
\label{fig:cg-signtest}
\end{figure}

\begin{happeningbox}
Three tests on the same eight differences:
paired $t$-test $p = 0.0026$; signed-rank $p = 0.016$; sign test $p = 0.070$.
Each discards more information than the one before. The $t$-test uses the
values, the signed-rank test uses only their order, and the sign test uses only
their signs. Less information means less power. The trade is robustness: an
outlier of $+500$ ms would wreck the $t$-test and leave the other two untouched.
Use the most informative test whose assumptions you believe.
\end{happeningbox}

The rank-based cousin of one-way ANOVA is the \textbf{Kruskal--Wallis} test.
It ranks all $N$ values together and runs the ANOVA idea on the ranks: its
statistic $H$ measures how far each group's average rank sits from the overall
average rank $(N+1)/2$, and under $H_0$ it is approximately $\chi^2_{k-1}$.
On the three server configurations it gives $H = 8.26$, $p = 0.016$, compared
with ANOVA's $p = 0.0047$.

\begin{lstlisting}[language=Python]
from scipy import stats
a = [12,15,19,22,30]; b = [18,25,28,35,40,41]
print(stats.mannwhitneyu(a, b, method='exact'))   # U = 5, p = 0.082
before = [120,135,128,150,142,138,160,131]
after  = [112,130,129,138,135,131,149,127]
print(stats.wilcoxon(before, after))              # p = 0.0156
print(stats.binomtest(7, 8))                      # sign test, p = 0.070
\end{lstlisting}

\section{A/B testing in practice}
\label{sec:cg-ab-practice}

The two-proportion test of Section~\ref{sec:cg-two-prop} is the arithmetic of
an A/B test. Getting a trustworthy answer takes more care than the arithmetic.

\subsection{The randomisation unit}

Every test so far assumes independent observations. In an A/B test, the
independent unit is whatever you randomised. If you assign by \emph{user}, the
same user stays in one arm for every visit, and users are the independent
units. If you assign by \emph{pageview}, one user can see both versions.

\begin{figure}[H]
\centering
\begin{tikzpicture}[pv/.style={minimum size=0.42cm,draw,inner sep=0pt},
  ua/.style={pv,fill=sourceblue!45},ub/.style={pv,fill=quotientorange!55}]
  \node[font=\small\bfseries] at (2.2,2.7) {randomise by user};
  \foreach \y/\s/\lab in {2/ua/1,1.3/ub/2,0.6/ua/3,-0.1/ub/4} {
    \node[font=\scriptsize,anchor=east] at (0.2,\y) {user \lab};
    \foreach \x in {0.6,1.1,1.6,2.1,2.6,3.1} \node[\s] at (\x,\y) {};
  }
  \begin{scope}[xshift=6.8cm]
  \node[font=\small\bfseries] at (2.2,2.7) {randomise by pageview};
  \foreach \y/\lab/\seq in {2/1/{0,1,1,0,1,0},1.3/2/{1,0,0,1,0,1},0.6/3/{0,0,1,1,0,1},-0.1/4/{1,1,0,0,1,0}} {
    \node[font=\scriptsize,anchor=east] at (0.2,\y) {user \lab};
    \foreach \c [count=\i] in \seq {
      \pgfmathsetmacro{\xx}{0.1+0.5*\i}
      \ifnum\c=0 \node[ua] at (\xx,\y) {}; \else \node[ub] at (\xx,\y) {}; \fi
    }
  }
  \end{scope}
  \node[pv,fill=sourceblue!45,label={[font=\scriptsize]right:A}] at (3.5,-0.9) {};
  \node[pv,fill=quotientorange!55,label={[font=\scriptsize]right:B}] at (5.0,-0.9) {};
\end{tikzpicture}
\caption{Left: each user sees one version on every visit; users are
independent units. Right: versions are mixed within a user, so one user's
experience of B can affect their behaviour on A, and their pageviews are not
independent.}
\label{fig:cg-unit}
\end{figure}

\begin{trapbox}
Analyse at the level you randomised. If you randomise users but compute the
standard error as if every pageview were independent, you pretend to have far
more data than you do. Pageviews from one user are correlated: a heavy user
contributes 200 similar pageviews. The SE comes out too small and false
positives multiply.
\end{trapbox}

\subsection{Ratio metrics and the delta method}

Many metrics are ratios of two per-user sums: clicks per session, revenue per
order, errors per request. Let user $i$ have $Y_i$ clicks and $X_i$ sessions.
The metric is
\[
  R = \frac{\sum_i Y_i}{\sum_i X_i} = \frac{\bar Y}{\bar X}.
\]
Both numerator and denominator are random, and they are correlated (users with
more sessions click more). The \emph{delta method} linearises the ratio around
the true means $\mu_Y, \mu_X$ with a first-order Taylor expansion, $r = \mu_Y/\mu_X$:
\[
  \frac{\bar Y}{\bar X} - r \approx \frac{1}{\mu_X}\left[(\bar Y - \mu_Y) - r(\bar X - \mu_X)\right].
\]
The right side is an average of the per-user quantities $(Y_i - rX_i)/\mu_X$,
so its variance is
\[
  \Var(R) \approx \frac{1}{n\,\mu_X^2}\left[\Var(Y) - 2r\Cov(X,Y) + r^2\Var(X)\right].
\]

\textbf{Numbers.} Suppose $n = 1000$ users per arm, $\mu_X = 4$ sessions,
$\mu_Y = 2$ clicks, so $r = 0.5$. Take $\SD(Y) = 3$, $\SD(X) = 5$, correlation
$0.8$, so $\Cov(X,Y) = 0.8 \times 3 \times 5 = 12$. Then
\[
  \Var(R) \approx \frac{9 - 2(0.5)(12) + 0.25(25)}{1000 \times 16}
  = \frac{3.25}{16000}, \qquad \SE(R) \approx 0.0143.
\]
The naive approach treats the $4000$ sessions as independent Bernoulli trials
with rate $0.5$: $\SE = \sqrt{0.25/4000} \approx 0.0079$. The naive SE is
too small by a factor of $1.8$. A test using it would declare far too many
false wins.

\begin{intuitionbox}
The delta method says: a smooth function of averages behaves like an average
of a linearised per-unit quantity. Compute that per-user quantity, take its
variance, divide by $n$. It is the same $\sigma/\sqrt n$ logic as Chapter 6,
applied after a Taylor step. When in doubt, the bootstrap over users
(Chapter 8) gives the same answer with no algebra.
\end{intuitionbox}

\subsection{Novelty and primacy effects}

Users react to change itself. A new button colour gets curious clicks in week
one and fades to its true effect later (novelty). Or users dislike any change
at first and warm up (primacy). Either way, the first days mislead.

\begin{figure}[H]
\centering
\begin{tikzpicture}
\begin{axis}[width=11cm,height=5cm,xmin=0,xmax=29,ymin=0,ymax=4.6,
  axis lines=left,xlabel={day of experiment},ylabel={daily lift (\%)},
  xtick={1,7,14,21,28}]
  \addplot[only marks,mark=*,mark size=1.6pt,sourceblue,domain=1:28,samples=28]
    {1 + 3*exp(-(x-1)/4)};
  \addplot[thick,sourceblue,domain=1:28,samples=100] {1 + 3*exp(-(x-1)/4)};
  \addplot[dashed,densityred,thick,domain=0:29] {1};
  \node[anchor=south east,font=\small,densityred] at (axis cs:28,1.05) {long-run lift $1\%$};
  \node[anchor=west,font=\small,sourceblue] at (axis cs:2.5,3.6) {novelty: early users click out of curiosity};
\end{axis}
\end{tikzpicture}
\caption{An illustrative novelty effect. Daily lift starts at $4\%$ and decays
to its long-run value of $1\%$. A test stopped after three days would report an
average lift of about $3.4\%$, more than three times the true value.}
\label{fig:cg-novelty}
\end{figure}

Defences: run for at least one or two full weekly cycles; plot the daily lift
and check that it has stabilised; compare new users (who never saw the old
version) with returning users.

\subsection{Sample ratio mismatch}

An experiment designed as a 50/50 split ends with $50{,}600$ users in A and
$49{,}400$ in B. That looks close. Is it? Run a chi-square goodness-of-fit
test against the designed split, expected $50{,}000$ each:
\[
  X^2 = \frac{600^2}{50000} + \frac{(-600)^2}{50000} = 7.2 + 7.2 = 14.4,
  \qquad \text{df} = 1, \qquad p \approx 0.00015.
\]
A split this uneven happens about once in 7000 experiments by chance. This is a
\textbf{sample ratio mismatch} (SRM). It means something in the pipeline is
dropping or duplicating users in one arm: a redirect that loses slow clients,
a bot filter that reacts to the new page, a crash in B that prevents logging.

\begin{trapbox}
If the SRM check fails, do not interpret the metrics. The users who
disappeared were not removed at random, so the arms are no longer comparable
and the randomisation that justified a causal claim is broken. Find the bug
first. Run the SRM check on every experiment before looking at any metric; it
costs one line of code.
\end{trapbox}

\begin{recoverybox}
Suppose the same $1.2\%$ imbalance ($50.6\%$ vs $49.4\%$) appears in an
experiment with only $1000$ users. Is that an SRM? (No. $X^2 = 6^2/500 \times 2 =
0.144$, $p \approx 0.70$. SRM tests become sharper as $n$ grows, which is why
large experiments catch small pipeline bugs.)
\end{recoverybox}

\subsection{CUPED: variance reduction with pre-experiment data}

Much of the variation in a user's spending this month is explained by their
spending last month. That variation is noise for the experiment, since it
existed before the treatment. CUPED (Controlled-experiment Using Pre-Experiment
Data) removes it.

Let $Y$ be the experiment metric and $X$ the same metric in the weeks before the
experiment. Define the adjusted metric
\[
  Y^{\text{adj}} = Y - \theta\,(X - \bar X), \qquad \theta = \frac{\Cov(X,Y)}{\Var(X)}.
\]
Since $X$ was measured before randomisation, its mean is the same in both arms
in expectation, so the adjustment does not bias the treatment effect. But it
shrinks the variance:
\[
  \Var(Y^{\text{adj}}) = \Var(Y)\,(1 - \rho^2),
\]
where $\rho$ is the correlation between $X$ and $Y$. The choice of $\theta$ is
the regression slope of $Y$ on $X$ (Chapter 12), which is exactly the value that
minimises this variance.

\begin{figure}[H]
\centering
\begin{tikzpicture}
\begin{axis}[width=8.5cm,height=6.5cm,xmin=28,xmax=68,ymin=28,ymax=72,
  axis lines=left,xlabel={pre-period metric $X$},ylabel={experiment metric $Y$}]
  \addplot[only marks,mark=*,mark size=1.6pt,sourceblue!80] coordinates
  {(39.0,47.5) (42.7,49.7) (42.2,60.0) (52.7,65.5) (47.5,54.8) (51.3,69.1) (58.4,67.2) (58.6,56.4) (54.8,67.7) (45.5,55.8) (42.5,45.2) (41.9,58.0) (46.6,55.1) (49.5,56.8) (40.3,57.4) (38.7,36.4) (53.1,53.0) (31.5,44.0) (48.2,60.5) (54.3,58.2) (40.1,46.0) (38.9,44.7) (42.4,53.8) (56.5,69.4) (48.7,48.7) (31.3,32.1) (45.8,54.9) (60.1,64.4) (59.8,60.4) (56.3,68.8) (47.6,64.1) (31.6,43.1) (51.7,58.7) (48.2,62.8) (50.8,57.5) (65.4,60.8) (51.8,60.8) (52.8,55.5) (56.1,66.6) (47.4,54.6)};
  \addplot[thick,densityred,domain=29:67] {15.889 + 0.8376*x};
  \draw[residgray,thick] (axis cs:31.3,32.1) -- (axis cs:31.3,42.11);
  \draw[residgray,thick] (axis cs:51.3,69.1) -- (axis cs:51.3,58.86);
  \node[anchor=west,font=\small,densityred] at (axis cs:50,38) {slope $\theta = 0.84$};
\end{axis}
\end{tikzpicture}
\caption{Forty users' pre-period and experiment-period values ($\rho = 0.77$).
CUPED keeps only each user's vertical distance from the red line (two shown in
grey). Those residuals vary much less than $Y$ itself: in this sample the
variance drops from $79.0$ to $32.4$, a factor of $0.41 = 1 - 0.77^2$.}
\label{fig:cg-cuped-scatter}
\end{figure}

\begin{figure}[H]
\centering
\begin{tikzpicture}
\begin{axis}[width=10.5cm,height=5cm,domain=-1:3,samples=150,
  axis lines=left,xlabel={estimated lift},ylabel={density},ytick=\empty,
  ymin=0,ymax=1.9,clip=false,legend style={at={(0.98,0.98)},anchor=north east,font=\small}]
  \addplot[thick,residgray] {exp(-(x-1)^2/(2*0.25))/sqrt(2*pi*0.25)};
  \addlegendentry{raw $Y$: SE $=0.50$}
  \addplot[thick,targetgreen] {exp(-(x-1)^2/(2*0.1025))/sqrt(2*pi*0.1025)};
  \addlegendentry{CUPED: SE $=0.32$}
  \draw[dashed] (axis cs:0,0) -- (axis cs:0,1.3);
  \node[anchor=south,font=\scriptsize] at (axis cs:0,1.3) {no effect};
\end{axis}
\end{tikzpicture}
\caption{Sampling distributions of the estimated lift (true lift 1) without and
with CUPED, for $\rho = 0.77$. The SE shrinks by $\sqrt{0.41} \approx 0.64$. The
CUPED estimate is far less likely to land near zero, so the same experiment has
much more power.}
\label{fig:cg-cuped-density}
\end{figure}

\begin{whybox}
Why $1 - \rho^2$? $\Var(Y - \theta X) = \Var(Y) - 2\theta\Cov(X,Y) +
\theta^2\Var(X)$. This is a quadratic in $\theta$, minimised at
$\theta = \Cov(X,Y)/\Var(X)$. Plugging that in gives
$\Var(Y) - \Cov(X,Y)^2/\Var(X) = \Var(Y)(1-\rho^2)$. Since required sample size
scales with variance, $\rho = 0.7$ cuts the users needed by half
($1 - 0.49 = 0.51$). CUPED is regression adjustment; it is the ``control for a
pre-treatment covariate'' idea of Chapter 13 applied to experiments.
\end{whybox}

\section{Which test when}
\label{sec:cg-which}

\begin{figure}[H]
\centering
\begin{tikzpicture}[font=\scriptsize,
  q/.style={draw,rounded corners,fill=structurepurple!10,align=center,text width=#1,inner sep=3pt},
  q/.default=1.9cm,
  leaf/.style={draw,fill=targetgreen!12,align=left,text width=4.1cm,inner sep=3pt},
  ar/.style={-{Stealth},thick,gray!70!black}]
  \node[q=1.4cm] (root) at (0.8,0) {What is the outcome?};
  \node[q=1.7cm] (cat) at (3.3,2.6) {Categorical (counts)};
  \node[q=1.7cm] (num) at (3.3,-2.6) {Numeric};
  \node[q=2.1cm] (c1) at (5.9,4.4) {One variable\\vs.\ claimed shares};
  \node[q=2.1cm] (c2) at (5.9,2.6) {Two groups, yes/no outcome};
  \node[q=2.1cm] (c3) at (5.9,0.8) {Two categorical variables ($r\times c$)};
  \node[q=2.1cm] (n1) at (5.9,-0.8) {Paired / before--after};
  \node[q=2.1cm] (n2) at (5.9,-2.6) {Two\\independent groups};
  \node[q=2.1cm] (n3) at (5.9,-4.4) {Three or more groups};
  \node[leaf] (l1) at (10.3,4.4) {Chi-square goodness of fit};
  \node[leaf] (l2) at (10.3,2.6) {Two-proportion $z$ ($=2\times2$ $\chi^2$)\\Fisher if any $E<5$};
  \node[leaf] (l3) at (10.3,0.8) {Chi-square independence\\Fisher if any $E<5$};
  \node[leaf] (l4) at (10.3,-0.8) {Normal-ish $d$: paired $t$\\Skewed: signed-rank, sign test};
  \node[leaf] (l5) at (10.3,-2.6) {Normal-ish: Welch $t$\\Skewed: Mann--Whitney};
  \node[leaf] (l6) at (10.3,-4.4) {Normal-ish: ANOVA + Tukey\\Skewed: Kruskal--Wallis};
  \draw[ar] (root) -- (cat); \draw[ar] (root) -- (num);
  \foreach \c in {c1,c2,c3} \draw[ar] (cat.east) -- (\c.west);
  \foreach \c in {n1,n2,n3} \draw[ar] (num.east) -- (\c.west);
  \draw[ar] (c1) -- (l1); \draw[ar] (c2) -- (l2); \draw[ar] (c3) -- (l3);
  \draw[ar] (n1) -- (l4); \draw[ar] (n2) -- (l5); \draw[ar] (n3) -- (l6);
\end{tikzpicture}
\caption{Choosing a test. First ask what kind of outcome you have, then how the
data are structured, then whether normal-theory assumptions are plausible. In
every A/B setting, also check the randomisation unit and run the SRM test
first.}
\label{fig:cg-flowchart}
\end{figure}

\begin{trapbox}
Running several tests and reporting the one with the smallest p-value is
p-hacking. Choose the test before seeing the data, from the structure of the
design, using a chart like Figure~\ref{fig:cg-flowchart}. Chapter 11 covers
why this matters and how badly it inflates false positives.
\end{trapbox}

A p-value from any of these tests answers ``is there a difference?'' It does
not answer ``how big, and does it matter?'' Chapter 11 turns to effect sizes,
to what happens when you run many tests at once, and to the common ways a
correct test still produces a misleading conclusion.

\section{Summary}

\begin{center}
\begin{tabular}{p{3.3cm}p{4.6cm}p{4.6cm}}
\toprule
Test & Statistic & Null distribution \\
\midrule
Two proportions & $z = (\hat p_B - \hat p_A)/\SE_0$, pooled $\hat p$ & $\Normal(0,1)$ \\
Goodness of fit & $X^2 = \sum (O-E)^2/E$ & $\chi^2_{k-1}$ \\
Independence & same, $E_{ij} = R_iC_j/N$ & $\chi^2_{(r-1)(c-1)}$ \\
Fisher exact & count in one cell, margins fixed & hypergeometric \\
One-way ANOVA & $F = \text{MSB}/\text{MSW}$ & $F_{k-1,\,N-k}$ \\
Mann--Whitney & $U$ = pairs with A $>$ B & exact / normal \\
Signed-rank & $W_+$ = sum of positive ranks & exact / normal \\
Sign test & number of positive differences & $\Bin(n, 1/2)$ \\
\bottomrule
\end{tabular}
\end{center}

\begin{itemize}
  \item Every test compares observed data with what $H_0$ predicts, scaled by noise.
  \item Pooled SE for the test of two proportions, unpooled SE for the interval.
    The $2\times 2$ chi-square statistic equals $z^2$.
  \item Chi-square tests need counts and expected counts of at least about 5;
    otherwise use Fisher's exact test.
  \item ANOVA splits $\text{SST} = \text{SSB} + \text{SSW}$; $F$ compares
    between-group to within-group variance. With two groups, $F = t^2$.
    A significant $F$ needs post-hoc tests to say which groups differ.
  \item Rank tests trade some power for robustness to skew and outliers.
  \item In A/B tests: analyse at the randomisation unit, use the delta method
    (or bootstrap) for ratio metrics, watch for novelty, check SRM before
    anything else, and use CUPED to cut variance by $1-\rho^2$.
\end{itemize}

\section{Exercises}

\begin{enumerate}
  \item A sign-up page converts 200 of 2000 visitors in A and 250 of 2000 in B.
  Compute the pooled rate, the null SE, $z$, and the two-sided p-value.
  \item A coin lands heads 60 times in 100 flips. Compute the chi-square
  goodness-of-fit statistic against a fair coin and its p-value. Compare with
  the one-proportion $z$ statistic.
  \item A $2\times 2$ table has row totals 30 and 70, column totals 40 and 60,
  grand total 100. Find all four expected counts under independence. How many
  degrees of freedom does the test have?
  \item An ANOVA with $k = 4$ groups and $N = 24$ observations has
  $\text{SSB} = 24$ and $\text{SSW} = 60$. Complete the ANOVA table. Is the
  result significant at 5\%? (The 95th percentile of $F_{3,20}$ is $3.10$.)
  \item A two-group ANOVA gives $F = 6.25$. What is $|t|$ for the pooled
  two-sample $t$-test on the same data?
  \item Samples A: $3, 8, 10$ and B: $5, 12, 14, 20$. Compute the rank sum of A
  and $U_A$. What fraction of (A, B) pairs have A larger?
  \item In 10 paired comparisons, 9 differences are positive. Compute the
  two-sided sign test p-value.
  \item A 50/50 experiment ends with $10{,}300$ users in A and $9{,}700$ in B.
  Run the SRM check. What should you do next?
  \item Pre-period and experiment metrics have correlation $\rho = 0.5$. By
  what factor does CUPED reduce the variance? Repeat for $\rho = 0.9$.
  \item Why should you not run three pairwise $t$-tests at $\alpha = 0.05$ to
  compare three groups? Estimate the chance of at least one false positive if
  the tests were independent.
\end{enumerate}

\subsection*{Answers}

\begin{enumerate}
  \item $\hat p = 450/4000 = 0.1125$;
  $\SE_0 = \sqrt{0.1125 \times 0.8875 \times 2/2000} \approx 0.00999$;
  $z = 0.025/0.00999 \approx 2.50$; $p \approx 0.012$.
  \item $X^2 = 10^2/50 + 10^2/50 = 4$, df $= 1$, $p \approx 0.046$.
  $z = (0.6 - 0.5)/\sqrt{0.25/100} = 2$ and $z^2 = 4$.
  \item $E = 30\cdot40/100 = 12$, $30\cdot60/100 = 18$, $70\cdot40/100 = 28$,
  $70\cdot60/100 = 42$. df $= 1$.
  \item MSB $= 24/3 = 8$, MSW $= 60/20 = 3$, $F = 2.67 < 3.10$: not
  significant ($p \approx 0.076$).
  \item $|t| = \sqrt{6.25} = 2.5$.
  \item Pooled order: 3, 5, 8, 10, 12, 14, 20. A has ranks 1, 3, 4, so
  $R_A = 8$ and $U_A = 8 - 6 = 2$. A is larger in $2/12 \approx 17\%$ of pairs
  (8 > 5 and 10 > 5).
  \item $2(\binom{10}{9} + \binom{10}{10})/1024 = 22/1024 \approx 0.021$.
  \item $X^2 = 300^2/10000 \times 2 = 18$, $p \approx 0.00002$. SRM detected:
  do not read the metrics; find the logging or assignment bug first.
  \item $1 - 0.25 = 0.75$ (25\% less variance); $1 - 0.81 = 0.19$ (81\% less,
  about 5 times fewer users needed).
  \item The familywise false positive rate inflates: $1 - 0.95^3 \approx 0.14$.
  Use ANOVA, then a post-hoc procedure such as Tukey's HSD.
\end{enumerate}
